% FILE: CSE-22_TermFinal.tex
\documentclass{article}
\usepackage{amsmath}
\usepackage{amssymb}
\begin{document}

%%%%%%%%%%%%%%%%%%%%%%%%%%%%%%%%%%%%%%%%%%%%%%%%%%
% QUESTION_START
%%%%%%%%%%%%%%%%%%%%%%%%%%%%%%%%%%%%%%%%%%%%%%%%%%
% id=cse22-final-q5i
% discipline=CSE
% year=2022
% exam_type=Term Final
% section=B
% ct_number=UNKNOWN
% question_number=5i
% topic=Indefinite Integration
% subtopic=General
% difficulty=Easy
% length=Short
% question_type=Repeated PYQ Pattern
% frequency=1
% appearances=
% OCR_STATUS=VERIFIED
\section*{Term Final (Sec B) - Q5(i)}
Question:
Integrate the following:
\[
\int \frac{dx}{1+\cos x}
\]

Solution:
Step 1: Simplify the integrand using trigonometric half-angle identities.
Recall that \(1 + \cos x = 2 \cos^2\left(\frac{x}{2}\right)\).
Substituting this into the integral:
\[
\int \frac{dx}{1+\cos x} = \int \frac{dx}{2 \cos^2\left(\frac{x}{2}\right)}
\]

Step 2: Rewrite in terms of secant and integrate.
\[
\frac{1}{2} \int \sec^2\left(\frac{x}{2}\right) \, dx
\]
Using the standard integration formula \(\int \sec^2(u) \, du = \tan(u) + C\):
\[
\frac{1}{2} \cdot \left( \frac{\tan\left(\frac{x}{2}\right)}{\frac{1}{2}} \right) + C = \tan\left(\frac{x}{2}\right) + C
\]

Final Answer:
\[
\boxed{\tan\left(\frac{x}{2}\right) + C}
\]
%%%%%%%%%%%%%%%%%%%%%%%%%%%%%%%%%%%%%%%%%%%%%%%%%%
% QUESTION_END
%%%%%%%%%%%%%%%%%%%%%%%%%%%%%%%%%%%%%%%%%%%%%%%%%%

%%%%%%%%%%%%%%%%%%%%%%%%%%%%%%%%%%%%%%%%%%%%%%%%%%
% QUESTION_START
%%%%%%%%%%%%%%%%%%%%%%%%%%%%%%%%%%%%%%%%%%%%%%%%%%
% id=cse22-final-q5ii
% discipline=CSE
% year=2022
% exam_type=Term Final
% section=B
% ct_number=UNKNOWN
% question_number=5ii
% topic=Indefinite Integration
% subtopic=Substitution
% difficulty=Easy
% length=Short
% question_type=Repeated PYQ Pattern
% frequency=1
% appearances=
% OCR_STATUS=VERIFIED
\section*{Term Final (Sec B) - Q5(ii)}
Question:
Integrate the following:
\[
\int \frac{e^{2x}}{e^x+1} \, dx
\]

Solution:
Step 1: Apply substitution.
Let \(u = e^x\). Then \(du = e^x \, dx\).
We can rewrite the numerator as \(e^{2x} = e^x \cdot e^x\).
The integral becomes:
\[
\int \frac{e^x}{e^x+1} (e^x \, dx) = \int \frac{u}{u+1} \, du
\]

Step 2: Perform algebraic manipulation on the integrand.
\[
\int \frac{u+1-1}{u+1} \, du = \int \left( 1 - \frac{1}{u+1} \right) du
\]

Step 3: Integrate and substitute back \(u = e^x\).
\[
u - \ln|u+1| + C = e^x - \ln(e^x+1) + C
\]

Final Answer:
\[
\boxed{e^x - \ln(e^x+1) + C}
\]
%%%%%%%%%%%%%%%%%%%%%%%%%%%%%%%%%%%%%%%%%%%%%%%%%%
% QUESTION_END
%%%%%%%%%%%%%%%%%%%%%%%%%%%%%%%%%%%%%%%%%%%%%%%%%%

%%%%%%%%%%%%%%%%%%%%%%%%%%%%%%%%%%%%%%%%%%%%%%%%%%
% QUESTION_START
%%%%%%%%%%%%%%%%%%%%%%%%%%%%%%%%%%%%%%%%%%%%%%%%%%
% id=cse22-final-q5iii
% discipline=CSE
% year=2022
% exam_type=Term Final
% section=B
% ct_number=UNKNOWN
% question_number=5iii
% topic=Indefinite Integration
% subtopic=Integration by Parts
% difficulty=Medium
% length=Medium
% question_type=Repeated PYQ Pattern
% frequency=1
% appearances=
% OCR_STATUS=VERIFIED
\section*{Term Final (Sec B) - Q5(iii)}
Question:
Integrate the following:
\[
\int e^{2x} \sin 4x \, dx
\]

Solution:
Step 1: Use the standard formula for integrals of the type \(\int e^{ax} \sin bx \, dx\).
Recall the formula:
\[
\int e^{ax} \sin bx \, dx = \frac{e^{ax}}{a^2+b^2} (a \sin bx - b \cos bx) + C
\]

Step 2: Substitute the parameters \(a=2\) and \(b=4\) into the formula.
\[
\int e^{2x} \sin 4x \, dx = \frac{e^{2x}}{2^2+4^2} (2 \sin 4x - 4 \cos 4x) + C
\]
\[
= \frac{e^{2x}}{20} (2 \sin 4x - 4 \cos 4x) + C
\]

Step 3: Simplify the expression.
\[
= \frac{e^{2x}}{10} (\sin 4x - 2 \cos 4x) + C
\]

Final Answer:
\[
\boxed{\frac{e^{2x}}{10} (\sin 4x - 2 \cos 4x) + C}
\]
%%%%%%%%%%%%%%%%%%%%%%%%%%%%%%%%%%%%%%%%%%%%%%%%%%
% QUESTION_END
%%%%%%%%%%%%%%%%%%%%%%%%%%%%%%%%%%%%%%%%%%%%%%%%%%

%%%%%%%%%%%%%%%%%%%%%%%%%%%%%%%%%%%%%%%%%%%%%%%%%%
% QUESTION_START
%%%%%%%%%%%%%%%%%%%%%%%%%%%%%%%%%%%%%%%%%%%%%%%%%%
% id=cse22-final-q5iv
% discipline=CSE
% year=2022
% exam_type=Term Final
% section=B
% ct_number=UNKNOWN
% question_number=5iv
% topic=Indefinite Integration
% subtopic=Partial Fractions
% difficulty=Medium
% length=Medium
% question_type=Repeated PYQ Pattern
% frequency=1
% appearances=
% OCR_STATUS=VERIFIED
\section*{Term Final (Sec B) - Q5(iv)}
Question:
Integrate the following:
\[
\int \frac{x \, dx}{x^2-12x+35}
\]

Solution:
Step 1: Factor the denominator.
The quadratic expression \(x^2-12x+35\) can be factored as:
\[
x^2-12x+35 = (x-5)(x-7)
\]

Step 2: Apply partial fraction decomposition.
\[
\frac{x}{(x-5)(x-7)} = \frac{A}{x-5} + \frac{B}{x-7}
\]
Multiply both sides by \((x-5)(x-7)\):
\[
x = A(x-7) + B(x-5)
\]

Step 3: Solve for constants \(A\) and \(B\).
For \(x = 5\):
\[
5 = A(5-7) \implies 5 = -2A \implies A = -\frac{5}{2}
\]
For \(x = 7\):
\[
7 = B(7-5) \implies 7 = 2B \implies B = \frac{7}{2}
\]

Step 4: Integrate the decomposed terms.
\[
\int \frac{x \, dx}{x^2-12x+35} = -\frac{5}{2} \int \frac{dx}{x-5} + \frac{7}{2} \int \frac{dx}{x-7}
\]
\[
= -\frac{5}{2} \ln|x-5| + \frac{7}{2} \ln|x-7| + C
\]

Final Answer:
\[
\boxed{\frac{7}{2} \ln|x-7| - \frac{5}{2} \ln|x-5| + C}
\]
%%%%%%%%%%%%%%%%%%%%%%%%%%%%%%%%%%%%%%%%%%%%%%%%%%
% QUESTION_END
%%%%%%%%%%%%%%%%%%%%%%%%%%%%%%%%%%%%%%%%%%%%%%%%%%

%%%%%%%%%%%%%%%%%%%%%%%%%%%%%%%%%%%%%%%%%%%%%%%%%%
% QUESTION_START
%%%%%%%%%%%%%%%%%%%%%%%%%%%%%%%%%%%%%%%%%%%%%%%%%%
% id=cse22-final-q6a
% discipline=CSE
% year=2022
% exam_type=Term Final
% section=B
% ct_number=UNKNOWN
% question_number=6a
% topic=Definite Integration
% subtopic=General
% difficulty=Medium
% length=Medium
% question_type=Repeated PYQ Pattern
% frequency=1
% appearances=
% OCR_STATUS=VERIFIED
\section*{Term Final (Sec B) - Q6(a)}
Question:
Define definite integral. Find \(\int_1^2 x \, dx\) as the limit of sum.

Solution:
Step 1: State the definition of a definite integral.
Let \(f(x)\) be a continuous function defined on a closed interval \([a, b]\). The definite integral of \(f(x)\) from \(a\) to \(b\) is defined as the limit of Riemann sums as the width of the subintervals approaches zero:
\[
\int_a^b f(x) \, dx = \lim_{n \to \infty} \sum_{i=1}^{n} f(x_i) \Delta x
\]
where \(\Delta x = \frac{b-a}{n}\) and \(x_i = a + i\Delta x\).

Step 2: Set up the limit of sum for \(\int_1^2 x \, dx\).
Here, \(f(x) = x\), \(a = 1\), and \(b = 2\).
\[
\Delta x = \frac{2-1}{n} = \frac{1}{n}
\]
\[
x_i = 1 + i\Delta x = 1 + \frac{i}{n}
\]

Step 3: Substitute the variables into the summation formula.
\[
\int_1^2 x \, dx = \lim_{n \to \infty} \sum_{i=1}^{n} \left(1 + \frac{i}{n}\right) \frac{1}{n}
\]
\[
= \lim_{n \to \infty} \frac{1}{n} \left[ \sum_{i=1}^{n} 1 + \frac{1}{n} \sum_{i=1}^{n} i \right]
\]

Step 4: Evaluate the summations.
Recall that \(\sum_{i=1}^{n} 1 = n\) and \(\sum_{i=1}^{n} i = \frac{n(n+1)}{2}\).
\[
\int_1^2 x \, dx = \lim_{n \to \infty} \frac{1}{n} \left[ n + \frac{n(n+1)}{2n} \right]
\]
\[
= \lim_{n \to \infty} \left[ 1 + \frac{n+1}{2n} \right]
\]
\[
= \lim_{n \to \infty} \left[ 1 + \frac{1}{2} + \frac{1}{2n} \right] = 1 + \frac{1}{2} = \frac{3}{2}
\]

Final Answer:
\[
\boxed{\frac{3}{2}}
\]
%%%%%%%%%%%%%%%%%%%%%%%%%%%%%%%%%%%%%%%%%%%%%%%%%%
% QUESTION_END
%%%%%%%%%%%%%%%%%%%%%%%%%%%%%%%%%%%%%%%%%%%%%%%%%%

%%%%%%%%%%%%%%%%%%%%%%%%%%%%%%%%%%%%%%%%%%%%%%%%%%
% QUESTION_START
%%%%%%%%%%%%%%%%%%%%%%%%%%%%%%%%%%%%%%%%%%%%%%%%%%
% id=cse22-final-q6b
% discipline=CSE
% year=2022
% exam_type=Term Final
% section=B
% ct_number=UNKNOWN
% question_number=6b
% topic=Definite Integration
% subtopic=General
% difficulty=Hard
% length=Medium
% question_type=New
% frequency=1
% appearances=
% OCR_STATUS=VERIFIED
\section*{Term Final (Sec B) - Q6(b)}
Question:
Evaluate \(\lim_{n \to \infty} \left[\frac{1}{n+m} + \frac{1}{n+2m} + \frac{1}{n+3m} + \dots + \frac{1}{n+nm}\right]\).

Solution:
Step 1: Express the sequence as a summation.
The given expression can be written as:
\[
S_n = \sum_{r=1}^{n} \frac{1}{n+rm}
\]

Step 2: Convert the summation into a Riemann sum form.
Factor out \(n\) from the denominator:
\[
S_n = \sum_{r=1}^{n} \frac{1}{n\left(1 + m\frac{r}{n}\right)} = \frac{1}{n} \sum_{r=1}^{n} \frac{1}{1 + m\left(\frac{r}{n}\right)}
\]

Step 3: Convert the limit of the Riemann sum into a definite integral.
As \(n \to \infty\), let \(\frac{r}{n} \to x\) and \(\frac{1}{n} \to dx\). The limits of integration run from \(x = 0\) (for \(r=1\)) to \(x = 1\) (for \(r=n\)).
\[
\lim_{n \to \infty} S_n = \int_0^1 \frac{1}{1+mx} \, dx
\]

Step 4: Integrate and evaluate the expression.
Using the substitution method or standard integration rule:
\[
\int_0^1 \frac{1}{1+mx} \, dx = \left[ \frac{\ln(1+mx)}{m} \right]_0^1
\]
\[
= \frac{\ln(1+m)}{m} - \frac{\ln(1)}{m} = \frac{\ln(1+m)}{m}
\]

Final Answer:
\[
\boxed{\frac{\ln(1+m)}{m}}
\]
%%%%%%%%%%%%%%%%%%%%%%%%%%%%%%%%%%%%%%%%%%%%%%%%%%
% QUESTION_END
%%%%%%%%%%%%%%%%%%%%%%%%%%%%%%%%%%%%%%%%%%%%%%%%%%

%%%%%%%%%%%%%%%%%%%%%%%%%%%%%%%%%%%%%%%%%%%%%%%%%%
% QUESTION_START
%%%%%%%%%%%%%%%%%%%%%%%%%%%%%%%%%%%%%%%%%%%%%%%%%%
% id=cse22-final-q7a
% discipline=CSE
% year=2022
% exam_type=Term Final
% section=B
% ct_number=UNKNOWN
% question_number=7a
% topic=Definite Integration
% subtopic=General
% difficulty=Easy
% length=Medium
% question_type=Repeated PYQ Pattern
% frequency=1
% appearances=
% OCR_STATUS=VERIFIED
\section*{Term Final (Sec B) - Q7(a)}
Question:
What are beta function and gamma function? Write the relation between them.

Solution:
Step 1: Define the Beta Function.
The Beta Function, denoted by \(B(m, n)\), is a first-kind Eulerian integral defined for \(m > 0, n > 0\) as:
\[
B(m, n) = \int_0^1 x^{m-1} (1-x)^{n-1} \, dx
\]

Step 2: Define the Gamma Function.
The Gamma Function, denoted by \(\Gamma(n)\), is a second-kind Eulerian integral extension of the factorial function to real and complex numbers, defined for \(n > 0\) as:
\[
\Gamma(n) = \int_0^\infty e^{-x} x^{n-1} \, dx
\]

Step 3: Write the relationship between them.
The relationship connecting the Beta and Gamma functions is given by:
\[
B(m, n) = \frac{\Gamma(m) \Gamma(n)}{\Gamma(m+n)}
\]

Final Answer:
\[
\boxed{B(m, n) = \frac{\Gamma(m) \Gamma(n)}{\Gamma(m+n)}}
\]
%%%%%%%%%%%%%%%%%%%%%%%%%%%%%%%%%%%%%%%%%%%%%%%%%%
% QUESTION_END
%%%%%%%%%%%%%%%%%%%%%%%%%%%%%%%%%%%%%%%%%%%%%%%%%%

%%%%%%%%%%%%%%%%%%%%%%%%%%%%%%%%%%%%%%%%%%%%%%%%%%
% QUESTION_START
%%%%%%%%%%%%%%%%%%%%%%%%%%%%%%%%%%%%%%%%%%%%%%%%%%
% id=cse22-final-q7b
% discipline=CSE
% year=2022
% exam_type=Term Final
% section=B
% ct_number=UNKNOWN
% question_number=7b
% topic=Definite Integration
% subtopic=General
% difficulty=Hard
% length=Medium
% question_type=New
% frequency=1
% appearances=
% OCR_STATUS=VERIFIED
\section*{Term Final (Sec B) - Q7(b)}
Question:
Evaluate \(\int_0^1 \frac{dx}{(1-x^6)^{1/6}}\) by Beta-Gamma function.

Solution:
Step 1: Set up substitution to match the Beta function form.
Let \(x^6 = y \implies x = y^{1/6}\).
Differentiating both sides:
\[
dx = \frac{1}{6} y^{-5/6} \, dy
\]
When \(x=0, y=0\), and when \(x=1, y=1\).

Step 2: Express the integral in terms of \(y\).
\[
\int_0^1 \frac{dx}{(1-x^6)^{1/6}} = \int_0^1 \frac{\frac{1}{6} y^{-5/6}}{(1-y)^{1/6}} \, dy
\]
\[
= \frac{1}{6} \int_0^1 y^{-5/6} (1-y)^{-1/6} \, dy
\]

Step 3: Match with the standard Beta function definition.
Recall the definition: \(B(p, q) = \int_0^1 y^{p-1} (1-y)^{q-1} \, dy\).
Comparing the powers:
\[
p - 1 = -\frac{5}{6} \implies p = \frac{1}{6}
\]
\[
q - 1 = -\frac{1}{6} \implies q = \frac{5}{6}
\]
Thus, the integral is:
\[
\frac{1}{6} B\left(\frac{1}{6}, \frac{5}{6}\right) = \frac{1}{6} \frac{\Gamma(1/6)\Gamma(5/6)}{\Gamma(1/6+5/6)} = \frac{1}{6} \frac{\Gamma(1/6)\Gamma(5/6)}{\Gamma(1)}
\]

Step 4: Use Euler's reflection formula to evaluate the Gamma product.
Recall Euler's reflection formula: \(\Gamma(z)\Gamma(1-z) = \frac{\pi}{\sin(\pi z)}\).
Here, \(z = 1/6\):
\[
\Gamma(1/6)\Gamma(5/6) = \frac{\pi}{\sin(\pi/6)} = \frac{\pi}{1/2} = 2\pi
\]
Since \(\Gamma(1) = 1\), the integral evaluates to:
\[
\frac{1}{6} \cdot 2\pi = \frac{\pi}{3}
\]

Final Answer:
\[
\boxed{\frac{\pi}{3}}
\]
%%%%%%%%%%%%%%%%%%%%%%%%%%%%%%%%%%%%%%%%%%%%%%%%%%
% QUESTION_END
%%%%%%%%%%%%%%%%%%%%%%%%%%%%%%%%%%%%%%%%%%%%%%%%%%

%%%%%%%%%%%%%%%%%%%%%%%%%%%%%%%%%%%%%%%%%%%%%%%%%%
% QUESTION_START
%%%%%%%%%%%%%%%%%%%%%%%%%%%%%%%%%%%%%%%%%%%%%%%%%%
% id=cse22-final-q7c
% discipline=CSE
% year=2022
% exam_type=Term Final
% section=B
% ct_number=UNKNOWN
% question_number=7c
% topic=Applications of Derivatives
% subtopic=General
% difficulty=Medium
% length=Medium
% question_type=New
% frequency=1
% appearances=
% OCR_STATUS=VERIFIED
\section*{Term Final (Sec B) - Q7(c)}
Question:
Find the area of the region bounded above by \(y = x + b\), bounded below by \(y = x^2\), and bounded on the sides by the lines \(x = 0\) and \(x = 2\).

Solution:
Step 1: Set up the definite integral for the area.
The area \(A\) between two curves \(y_{top} = f(x)\) and \(y_{bottom} = g(x)\) on the interval \([a, b]\) is given by:
\[
A = \int_a^b [f(x) - g(x)] \, dx
\]
Substituting the given equations and boundaries:
\[
A = \int_0^2 [(x + b) - x^2] \, dx
\]

Step 2: Integrate the term with respect to \(x\).
\[
A = \left[ \frac{x^2}{2} + bx - \frac{x^3}{3} \right]_0^2
\]

Step 3: Evaluate the expression at the boundaries.
\[
A = \left( \frac{2^2}{2} + b(2) - \frac{2^3}{3} \right) - (0)
\]
\[
A = 2 + 2b - \frac{8}{3}
\]
\[
A = 2b - \frac{2}{3}
\]

Final Answer:
\[
\boxed{2b - \frac{2}{3}}
\]
%%%%%%%%%%%%%%%%%%%%%%%%%%%%%%%%%%%%%%%%%%%%%%%%%%
% QUESTION_END
%%%%%%%%%%%%%%%%%%%%%%%%%%%%%%%%%%%%%%%%%%%%%%%%%%

%%%%%%%%%%%%%%%%%%%%%%%%%%%%%%%%%%%%%%%%%%%%%%%%%%
% QUESTION_START
%%%%%%%%%%%%%%%%%%%%%%%%%%%%%%%%%%%%%%%%%%%%%%%%%%
% id=cse22-final-q8a
% discipline=CSE
% year=2022
% exam_type=Term Final
% section=B
% ct_number=UNKNOWN
% question_number=8a
% topic=Applications of Derivatives
% subtopic=General
% difficulty=Easy
% length=Medium
% question_type=Repeated PYQ Pattern
% frequency=1
% appearances=
% OCR_STATUS=VERIFIED
\section*{Term Final (Sec B) - Q8(a)}
Question:
Derive the formula for the volume of a sphere of radius \(r\).

Solution:
Step 1: Set up the geometric model.
Consider a circle with its center at the origin and radius \(r\), given by the equation:
\[
x^2 + y^2 = r^2 \implies y^2 = r^2 - x^2
\]
If we rotate the upper semi-circle \(y = \sqrt{r^2 - x^2}\) about the \(x\)-axis, it generates a solid sphere of radius \(r\).

Step 2: Formulate the volume using the disc method.
The volume of a representative disc of thickness \(dx\) and radius \(y\) is \(dV = \pi y^2 \, dx\).
Integrating this from \(x = -r\) to \(x = r\):
\[
V = \int_{-r}^{r} \pi y^2 \, dx = \pi \int_{-r}^{r} (r^2 - x^2) \, dx
\]

Step 3: Use symmetry and integrate.
Since the integrand is an even function, we can multiply by 2 and integrate from \(0\) to \(r\):
\[
V = 2\pi \int_0^r (r^2 - x^2) \, dx
\]
\[
V = 2\pi \left[ r^2 x - \frac{x^3}{3} \right]_0^r
\]

Step 4: Evaluate the integrated expression.
\[
V = 2\pi \left( r^3 - \frac{r^3}{3} \right)
\]
\[
V = 2\pi \left( \frac{2r^3}{3} \right) = \frac{4}{3}\pi r^3
\]

Final Answer:
\[
\boxed{\frac{4}{3}\pi r^3}
\]
%%%%%%%%%%%%%%%%%%%%%%%%%%%%%%%%%%%%%%%%%%%%%%%%%%
% QUESTION_END
%%%%%%%%%%%%%%%%%%%%%%%%%%%%%%%%%%%%%%%%%%%%%%%%%%

%%%%%%%%%%%%%%%%%%%%%%%%%%%%%%%%%%%%%%%%%%%%%%%%%%
% QUESTION_START
%%%%%%%%%%%%%%%%%%%%%%%%%%%%%%%%%%%%%%%%%%%%%%%%%%
% id=cse22-final-q8b
% discipline=CSE
% year=2022
% exam_type=Term Final
% section=B
% ct_number=UNKNOWN
% question_number=8b
% topic=Applications of Derivatives
% subtopic=General
% difficulty=Medium
% length=Medium
% question_type=Repeated PYQ Pattern
% frequency=1
% appearances=
% OCR_STATUS=VERIFIED
\section*{Term Final (Sec B) - Q8(b)}
Question:
Find the area of the surface that is generated by revolving the portion of the curve \(y = x^3\) between \(x = 0\) and \(x = 1\) about the \(x\)-axis.

Solution:
Step 1: State the formula for the surface area of revolution.
The surface area \(S\) of a solid generated by rotating a curve \(y = f(x)\) on the interval \([a, b]\) about the \(x\)-axis is given by:
\[
S = 2\pi \int_a^b y \sqrt{1 + \left(\frac{dy}{dx}\right)^2} \, dx
\]

Step 2: Compute the derivative of the given function.
Here, \(y = x^3 \implies \frac{dy}{dx} = 3x^2\).
Now, compute \(1 + \left(\frac{dy}{dx}\right)^2\):
\[
1 + (3x^2)^2 = 1 + 9x^4
\]

Step 3: Set up the integration.
Substitute these values into the surface area formula:
\[
S = 2\pi \int_0^1 x^3 \sqrt{1 + 9x^4} \, dx
\]

Step 4: Solve the integral using \(u\)-substitution.
Let \(u = 1 + 9x^4\).
Then \(du = 36x^3 \, dx \implies x^3 \, dx = \frac{du}{36}\).
Change the limits of integration:
When \(x = 0 \implies u = 1\).
When \(x = 1 \implies u = 10\).

Substituting these into the integral:
\[
S = 2\pi \int_1^{10} \sqrt{u} \frac{du}{36} = \frac{\pi}{18} \int_1^{10} u^{1/2} \, du
\]
\[
= \frac{\pi}{18} \left[ \frac{2}{3} u^{3/2} \right]_1^{10}
\]
\[
= \frac{\pi}{27} (10^{3/2} - 1^{3/2}) = \frac{\pi}{27} (10\sqrt{10} - 1)
\]

Final Answer:
\[
\boxed{\frac{\pi}{27} (10\sqrt{10} - 1)}
\]
%%%%%%%%%%%%%%%%%%%%%%%%%%%%%%%%%%%%%%%%%%%%%%%%%%
% QUESTION_END
%%%%%%%%%%%%%%%%%%%%%%%%%%%%%%%%%%%%%%%%%%%%%%%%%%

\end{document}